\clearpage
\appendix
\setcounter{page}{1}
% \maketitlesupplementary
\onecolumn

\begin{center}
    \Large
    \textbf{\thetitle} \\ % Use the same title as before or a new one
    \vspace{0.5em} % Space between title and subtitle
    Supplementary Material \\
    \vspace{1.0em} % Space after subtitle
\end{center}

In this supplementary file, we provide the following materials:

\begin{itemize}[leftmargin=2em]
    \item The interface and more details of the user study (referring to Sec. \ref{sec:4-1} in the main paper);
    \item Details of the evaluation metrics for the discriminator (referring to Sec. \ref{sec:4-1} and Sec. \ref{sec:train_discriminator} in the main paper);
    \item Visual examples of object removal results at each training round (referring to Sec \ref{sec:4-2} in the main paper);
    \item Visual results of IP2P and CLIPAway (referring to Sec \ref{sec:4-2} and Fig. \ref{fig:qualitative} in the main paper);
    \item Visual results on images of 512 × 512 resolution (referring to Sec \ref{sec:4-2} in the main paper);
    \item Failure cases (referring to Sec. \ref{sec:conclusion} in the main paper).
\end{itemize}

%%%%%%%%% BODY TEXT - ENTER YOUR RESPONSE BELOW
\section{Annotation page and user study page}
\label{sec:SP1}

We design a webpage based on the open-source library Gradio \cite{abid2019gradio} to conduct the human annotation (referring to Sec. \ref{sec:human_annotation} in the main paper) and the final human evaluation (referring to Sec. \ref{sec:4-1} in the main paper). Annotators are provided with the original images, the mask images and the object removal results, as illustrated in \cref{fig:humanAnnotation}. They are asked to provide feedback by clicking the \textbf{Yes} or \textbf{No} button at the bottom right corner.

\begin{figure*}[!h]
  \centering
   \includegraphics[width=\linewidth]{figures_appendix/Human_annotation_page.pdf}
\vspace{-3mm}
   \caption{The interface for human annotation. The annotators are asked to give feedback by clicking "Yes" or "No" button.}
   \label{fig:humanAnnotation}  
\end{figure*}

The interface for final human evaluation is shown in \cref{fig:humanEvaluation}. The input images and the masked images are displayed in the left column. The editing results of different methods as displayed in the right columns. Annotators are asked to click the multiple-choice check-boxes to select the successful removal results among different methods and submit the results. We randomly shuffle the display order in each evaluation. Five volunteers participated in the final evaluation, and each volunteer annotated 1,000 samples, including 500 pairs of object removal cases under $512 \times 512$ resolution and 500 pairs under $1024 \times 1024$ resolution. We calculate the average success rate for different methods based on these human evaluations.

\begin{figure*}[!htbp]
  \centering
   \includegraphics[width=\linewidth]{figures_appendix/Final_evaluation.pdf}
\vspace{-3mm}
   \caption{The interface for human evaluation. The volunteers make selections by checking the multiple-choice check-boxes at the bottom left corner.}
   \label{fig:humanEvaluation}  
  % \vspace{-5mm}
\end{figure*}

\section{The evaluation metrics for the discriminator}
\label{sec:SP2}

\begin{figure*}[!htbp]
  \centering
   \includegraphics[width=0.62\linewidth]{figures_appendix/discriminator_metrics.pdf}
   \caption{Confusion matrix and the definition of metrics for evaluating our discriminator. }
   \label{fig:evaluate_metrics}  
\end{figure*}

We use the 500 pairs in the test set with $512 \times 512$ resolution to test the discriminator. The edited results are generated by our RORem. The definitions of precision, recall, F1 and accuracy are illustrated in \cref{fig:evaluate_metrics}. Among these metrics, precision represents the percentage of the true positive samples to the total positive samples predicted by our discriminator. High precision ensures that the selected removal pairs are all of high-quality. By setting the threshold as 0.9, our final discriminator can reach a precision of 0.983, which allows us to obtain a large amount of high-quality data pairs. 

\section{Visual examples of object removal results at each training round}
\label{sec:SP3}

\begin{figure*}[thbp]
  \centering
   \includegraphics[width=\linewidth]{figures_appendix/step_result_final.pdf}
   \caption{Visual results of RORem at each training round, one can see that the removal capacity of RORem improves with the increase of the dataset. }
   \label{fig:step_result}  
\end{figure*}

The visual examples of object removal results at each training round are provided in \cref{fig:step_result}. We can see that the initial model SDXL-inpainting \cite{podell2023sdxl} always fills the masked regions with semantically similar contents instead of removing it (column Initialization). After the first round of finetuning, RORem can successfully remove the selected sofa (row 2) and cat (row 4); however, its removal capacity is not good enough, leading to failures in other cases (see partial removal cases bottle, statue and blurry synthesis case airplane). After we  extend the training dataset and conduct more finetuning rounds (see column R1-Human to column R3-Auto), RORem can successfully remove the masked regions in most cases. Finally, we collect images with 2K resolution from DIV2K \cite{Agustsson_2017_CVPR_Workshops} and Flicker2K \cite{timofte2017ntire} to conduct the final stage finetuning, where the removal capacity of RORem can be well preserved (see column Final) and the image quality can be improved (see Tab. \ref{tab:train_data} in the main paper).

\section{Visual results of IP2P and CLIPAway}
\label{sec:SP4}

The visual editing results of IP2P and CLIPAway on images of resolution of $1024 \times 1024$ are illustrated in \cref{fig:IP2P_CLIPAway}. We can see that IP2P fails in all cases and even changes the overall style of the given images (see images plate in column 1 and car in column 5). CLIPAway exhibits the same problem as PPT, which often fills the masked regions with incorrect contents (see images sofa, dog and car).

\begin{figure*}[thbp]
  \centering
   \includegraphics[width=\linewidth]{figures_appendix/appendix_clipaway.pdf}
   \caption{Visual results of IP2P and CLIPAway on $1024 \times 1024$ resolution images. }
   \label{fig:IP2P_CLIPAway}  
\end{figure*}

\section{Visual results on images of $512 \times 512$ resolution}
\label{sec:SP5}

The qualitative comparisons on images of $512 \times 512$ resolution are illustrated in \cref{fig:visual_512}. We can see that Lama can generate blurry synthesis outputs in some cases (see images koala in column 4 and plate in column 6). SDXL-INP, IP2P and INST fail in most cases. Moreover, as INST and IP2P are text-driven removal methods, the ambiguity of text instructions can lead to removal failures of selected objects (see images hot air ballon in column 3 and cup cake in column 6). IP2P not only fails to remove the select objects but also changes the overall style and details of the original images (see images hot air ballon, koala, and cup cake). PPT and CLIPAway can fill the masked regions with nonexistent contents in images bird (column 1), koala (column 4) and statue (column 5). DesignEdit succeeds in the first two removal cases, however it suffers from visual artifacts (see images koala, plate). In contrast, RORem successfully removes the selected objects in most cases. Meanwhile, our distilled RORem-4S model also works well in these cases with less time overhead.

\begin{figure*}[thbp]
  \centering
   \includegraphics[width=\linewidth]{figures_appendix/appendix_visual_512.pdf}
   \vspace{-10mm}
   \caption{Visual comparison of the object removal results by RORem and other methods on $512 \times 512$ resolution images. One can see there can be incomplete removal regions, blurry synthesis output, wrong removal target, and incorrect synthese contents in previous methods, while RORem demonstrate robust removal performance.}
   \label{fig:visual_512}  
\end{figure*}

\section{Failure Cases}
\label{sec:SP6}

As we stated in the conclusion section of the main paper, although RORem achieves great improvement on the overall removal performance, it may fail in cases when the background contains human fingers and faces. Some failure cases are depicted in \cref{fig:visual_failure}. Future work will be conducted for further improving the performance of RORem on these editing scenarios.

\clearpage

\begin{figure*}[thbp]
  \centering
   \includegraphics[width=\linewidth]{figures_appendix/appendix_failure.pdf}
   \caption{Failure cases of RORem.}
   \label{fig:visual_failure}  
\end{figure*}